% Fragmento de inserción breve. Desarrollo y procedencia en la misma carpeta.
\subsubsection{Mémoire arithmétique du canal multiplicatif}
\label{prd-arit:seccion}

La conservation de la signature régionale se prolonge à l'évaluation
arithmétique complète d'APP. Pour une position orientée
\(x=(i,j,d)\in\{1,\ldots,9\}^2\times\{N,E,S,O\}\), soit
\(\chi(x)=(a,b,\epsilon)\), où \(a\) est la coordonnée fixe,
\(b\) la coordonnée mobile et \(\epsilon\) son orientation :
\[
 \chi(i,j,N)=(j,i,-1),\quad \chi(i,j,S)=(j,i,+1),\qquad
 \chi(i,j,E)=(i,j,+1),\quad \chi(i,j,O)=(i,j,-1).
\]
L'étiquette d'axe \(\iota\in\{i,j\}\) conserve laquelle des
deux coordonnées originales se déplace. En écrivant
\(\langle b\rangle_9=1+((b-1)\bmod9)\), les opérateurs et
l'évaluation sont
\begin{align}
 \overline P(a,b,\epsilon)&=(a,\langle b+\epsilon\rangle_9,\epsilon),
 &\overline H(a,b,\epsilon)&=(a,b,-\epsilon),\label{prd-arit:PH}\\
 n(a,b)&=ab,&r(a,b)&=\langle ab\rangle_9,\qquad
 q(a,b)=\frac{ab-r(a,b)}9.\label{prd-arit:rq}
\end{align}

\begin{proposition}[Représentation exacte du transport arithmétique]
\label{prd-arit:minimal}
On a \(\chi P=\overline P\chi\) et
\(\chi H=\overline H\chi\). La représentation \(\chi\)
possède \(162\) états et est la moins fine des représentations
stables sous \(P,H\) qui conservent le produit entier ou,
de manière équivalente, le couple \((r,q)\). Conserver également \(\iota\)
permet de retrouver les \(324\) positions orientées originales.
\end{proposition}
\begin{proof}
L'avancée modifie seulement \(b\), et le demi-tour seulement
\(\epsilon\), ce qui démontre les deux identités de transport.
Les neuf observations \(n_k=n(\overline P^k\chi(x))\),
\(0\leq k<9\), parcourent \(a,2a,\ldots,9a\). On retrouve ainsi
\(a=\min_k n_k\), \(b=n_0/a\) et le signe par
\((n_1/a-n_0/a)\bmod9\in\{1,8\}\). Toute équivalence stable
conservant \(n\) conserve ces neuf observations et distingue donc
tous les états \((a,b,\epsilon)\).
L'identité \(n=r+9q\) démontre l'équivalence des deux
évaluations. Chaque fibre de \(\chi\) contient les deux chemins
transposés : \((i,j,d)\mapsto(j,i,d^\top)\), avec
\(N^\top=O,O^\top=N,E^\top=S,S^\top=E\). L'étiquette d'axe distingue
ces deux préimages et permet de reconstruire la position et la direction.
\end{proof}

Le sélecteur TRIT active le produit en \(t\equiv2,5,8\pmod9\),
avant chaque avancée. Le demi-tour agit à l'achèvement de chaque
intervalle de \(27\) instants. Si \(y_{t-1}\) désigne l'état arithmétique
précédent, l'incrément et sa mémoire accumulée sont déterminés par
\begin{equation}
 \eta_q(t)=\mathbf1_{\{t\equiv2\ (3)\}}q(y_{t-1}),\qquad
 \lambda_q(t)=\lambda_q(t-1)+\eta_q(t),\quad \lambda_q(0)=0.
 \label{prd-arit:incremento}
\end{equation}
On a \(0\leq q(a,b)\leq a-1\). Pour inverser une
mise à jour, on annule d'abord le demi-tour correspondant,
puis l'avancée, et l'on soustrait enfin le quotient retrouvé de
l'état précédent. Le registre conserve la phase, l'axe, l'ordre des événements et
la mémoire ; \(\lambda_q\) est un cocycle additif de ce transport.

Chaque bloc de \(27\) contient neuf avancées et parcourt une fois
la coordonnée mobile. Le retour de position et d'orientation
en \(54\), et de phase en \(108\), coexiste donc avec
\begin{equation}
 \lambda_q(108N)=NQ_\times^+(a),\qquad
 Q_\times^+(a)=\frac{180a-4\sum_{b=1}^9r(a,b)}9,
 \qquad N\geq0.
 \label{prd-arit:retorno}
\end{equation}
En effet, chaque bloc somme les produits \(a,2a,\ldots,9a\) ;
sommer \(n=r+9q\) sur les quatre blocs donne la formule. Si
\(a\) est premier avec \(9\), les résidus ont pour somme \(45\), et
\(Q_\times^+(a)=20(a-1)\). Dans la fenêtre \(j\), comptée à partir de
\(j=0\), le signe relatif du calendrier est
\(\sigma_j=(-1)^{\lfloor j/3\rfloor}\), et l'orientation
effective est \(\epsilon_j=\epsilon_0\sigma_j\).
La charge signée par \(\sigma_j\) s'annule ; la charge pondérée
par l'orientation effective s'annule également lorsqu'elle est multipliée par
\(\epsilon_0\). Le registre conserve l'orientation initiale, la charge
accumulée et ses incréments. Ces quantités appartiennent
au curseur multiplicatif ; les lecteurs d'une trajectoire qui
alterne les deux feuillets agissent sur le registre de cette trajectoire.
